\begin{table}[H]
    \centering
    \caption{Desempenho dos modelos no conjunto de teste (2022), com variáveis individuais e municipais.}
    \label{tab-res-desempenho}
    \footnotesize
    \renewcommand{\arraystretch}{1.2}
    \setlength{\tabcolsep}{4pt}
    \begin{tabular}{l c c c c c c c}
        \hline
        \textbf{Modelo} & \textbf{AUC (IC 95\%)} & \textbf{p} & \textbf{Sens.} & \textbf{Espec.} & \textbf{Prec.} & \textbf{F1} & \textbf{Brier} \\
        \hline
        Regressão Logística & 0,707 (0,700–0,713) & < 0,001 & 0,623 & 0,683 & 0,301 & 0,406 & 0,214 \\
        Árvore de Decisão & 0,701 (0,694–0,707) & < 0,001 & 0,615 & 0,678 & 0,295 & 0,399 & 0,216 \\
        Random Forest & 0,716 (0,709–0,723) & < 0,001 & 0,637 & 0,676 & 0,302 & 0,409 & 0,210 \\
        MLP & 0,713 (0,706–0,720) & < 0,001 & 0,665 & 0,648 & 0,293 & 0,407 & 0,219 \\
        LightGBM & 0,724 (0,718–0,731) & -- & 0,640 & 0,688 & 0,311 & 0,418 & 0,202 \\
        \hline
    \end{tabular}
    \legend{Nota: métricas calculadas com limiar de 0,5. IC 95\% da AUC pelo método de DeLong; p = comparação da AUC com a do LightGBM (teste de DeLong). Sens. = sensibilidade; Espec. = especificidade; Prec. = precisão; Brier = escore de Brier (quanto menor, melhor).\\Fonte: elaborado pelo autor.}
\end{table}
